\documentclass[aspectratio=169,10pt]{beamer}

% ============================================================
%  PACKAGES
% ============================================================
\usepackage[T1]{fontenc}
\usepackage[utf8]{inputenc}
\usepackage{amsmath,amssymb,bm}
\usepackage{tikz}
\usetikzlibrary{arrows.meta,positioning,calc,backgrounds}
\usepackage{tcolorbox}
\tcbuselibrary{skins,breakable}
\usepackage{xcolor}

% ============================================================
%  COLORS
% ============================================================
\definecolor{darkbg}{HTML}{1A1A2E}
\definecolor{accent}{HTML}{E94560}
\definecolor{accentteal}{HTML}{0F9B8E}
\definecolor{accentamber}{HTML}{F5A623}
\definecolor{accentviolet}{HTML}{7B68EE}
\definecolor{textlight}{HTML}{EAEAEA}
\definecolor{textgray}{HTML}{9E9E9E}
\definecolor{lightbg}{HTML}{F7F7F7}
\definecolor{coralbox}{HTML}{FF6B6B}

% ============================================================
%  BEAMER SETUP
% ============================================================
\setbeamercolor{background canvas}{bg=lightbg}
\setbeamercolor{normal text}{fg=darkbg}
\setbeamertemplate{navigation symbols}{}
\setbeamertemplate{footline}{%
  \hfill\textcolor{textgray}{\tiny\texttt{mbastidaso@unal.edu.co}\quad\insertframenumber/\inserttotalframenumber\quad}%
  \vspace{4pt}
}
\setbeamertemplate{itemize item}{\textcolor{accent}{\textbullet}}
\setbeamertemplate{itemize subitem}{\textcolor{accentteal}{--}}

% ============================================================
%  DARK SLIDE ENVIRONMENT (decorative circles)
% ============================================================
\newenvironment{darkslide}{%
  \setbeamercolor{background canvas}{bg=darkbg}
  \setbeamercolor{normal text}{fg=textlight}
  \usebeamercolor[fg]{normal text}
  \begin{frame}
  \begin{tikzpicture}[remember picture,overlay]
    \fill[accent,opacity=0.10]   (current page.north east) ++(-1.2,-1.0) circle (1.7);
    \fill[accentteal,opacity=0.10](current page.south west) ++(1.6,1.2) circle (2.3);
    \fill[accentviolet,opacity=0.10](current page.south east) ++(-2.0,2.4) circle (1.1);
  \end{tikzpicture}
}{%
  \end{frame}
}

% ============================================================
%  CONTENT BOXES
% ============================================================
\newtcolorbox{defbox}[1][]{colback=accentviolet!8, colframe=accentviolet,
  boxrule=1pt, arc=3pt, left=6pt,right=6pt,top=5pt,bottom=5pt, #1}
\newtcolorbox{resbox}[1][]{colback=accentteal!8, colframe=accentteal,
  boxrule=1pt, arc=3pt, left=6pt,right=6pt,top=5pt,bottom=5pt, #1}
\newtcolorbox{obsbox}[1][]{colback=accentamber!10, colframe=accentamber,
  boxrule=1pt, arc=3pt, left=6pt,right=6pt,top=5pt,bottom=5pt, #1}
\newtcolorbox{taskbox}[1][]{colback=coralbox!10, colframe=coralbox,
  boxrule=1pt, arc=3pt, left=6pt,right=6pt,top=5pt,bottom=5pt, #1}

\newcommand{\R}{\mathbb{R}}

% ============================================================
\begin{document}

% ---------- TITLE ----------
\begin{darkslide}
  \vfill
  \begin{flushleft}
  {\Large\textcolor{accent}{\textbf{Teorema de Aproximaci\'on Universal}}}\\[6pt]
  {\large\textcolor{textlight}{\textbf{?`Qu\'e puede representar una red neuronal?}}}\\[14pt]
  {\small\textcolor{textgray}{Programaci\'on Cient\'ifica 2026-1 \quad$\cdot$\quad Universidad Nacional de Colombia}}
  \end{flushleft}
  \vfill
\end{darkslide}

% ---------- SLIDE 1: el hilo ----------
\begin{frame}{El hilo del curso}
  Toda nuestra aproximaci\'on tiene la misma forma:
  \[
    \hat{f}(x) = \sum_{i=1}^{n} c_i\,\phi_i(x).
  \]
  \medskip
  \begin{itemize}
    \item Con \textbf{bases fijas} (monomios, splines, B\'ezier, Fourier) elegimos los $\phi_i$ de antemano y solo ajustamos los $c_i$.
    \item En una \textbf{red neuronal} los $\phi_i(x)=\sigma(w_i^\top x + b_i)$ \emph{tambi\'en} se aprenden.
  \end{itemize}
  \medskip
  \begin{obsbox}
  \textbf{La pregunta de hoy:} si dejamos que la red elija sus propias bases, ?`qu\'e funciones puede aproximar? ?`Hay un l\'imite?
  \end{obsbox}
\end{frame}

% ---------- SLIDE 2: enunciado ----------
\begin{frame}{El teorema (Cybenko 1989, Hornik 1991)}
  Una red de \textbf{una sola capa oculta} con $n$ neuronas representa
  \[
    \hat{f}(x) = \sum_{i=1}^{n} c_i\,\sigma(w_i^\top x + b_i).
  \]
  \begin{defbox}
  \textbf{Teorema.} Sea $\sigma$ continua, acotada y no constante, y $K\subset\R^d$ compacto. Para toda funci\'on continua $f:K\to\R$ y todo $\varepsilon>0$, existe $n$ y par\'ametros $\{c_i,w_i,b_i\}$ tales que
  \[
    \sup_{x\in K}\bigl|\,f(x)-\hat{f}(x)\,\bigr| < \varepsilon.
  \]
  \end{defbox}
  \begin{obsbox}
  El conjunto de redes de una capa es \textbf{denso} en $C(K)$ con la norma del supremo. No se acota $n$: la garant\'ia es de \emph{existencia}, no de eficiencia.
  \end{obsbox}
\end{frame}


% ---------- SLIDE 4: conexion con la clase de redes ----------
\begin{frame}{Conexi\'on con la clase}
  \begin{columns}[T]
    \column{0.5\textwidth}
    \begin{resbox}
    \textbf{Lo que garantiza el teorema}\\[2pt]
    Existe un $n$ con error $<\varepsilon$
    \end{resbox}
    \column{0.5\textwidth}
    \begin{taskbox}
    \textbf{Lo que el teorema \emph{no} dice}\\[2pt]
    No dice cu\'antas neuronas, ni que el descenso de gradiente \emph{encuentre} esos par\'ametros. Existencia $\neq$ optimizaci\'on.
    \end{taskbox}
  \end{columns}
  \medskip
  \begin{obsbox}
  En la pr\'actica la profundidad importa: redes profundas logran el mismo $\varepsilon$ con muchas menos neuronas que una sola capa ancha.
  \end{obsbox}
\end{frame}



\end{document}
